%---------------------------------------------------------------------
% Part two: the Madani period. Seven questions.
%---------------------------------------------------------------------

\section*{حصہ دوم --- مدنی دور}
\markboth{مدنی دور}{}

%---------------------------------------------------------------------
\begin{sawal}{سوال \textenglish{6} \ \textnormal{\small(بہار \textenglish{2022}، سوال \textenglish{2} اور خزاں \textenglish{2017}، سوال \textenglish{1}(الف\textenglish{3}))}}
ہجرتِ مدینہ کا واقعہ اور بیعتِ عقبہ پر جامع نوٹ لکھیں۔
\end{sawal}

\begin{hal}
\textbf{بیعتِ عقبہ --- دو مرحلے، ہجرت سے پہلے:}
\begin{itemize}\setlength{\itemsep}{1pt}
  \item \textbf{عقبہ اولیٰ} (\textenglish{621}ء) --- یثرب کے بارہ افراد نے
        شرک، چوری، زنا اور اولاد کے قتل سے باز رہنے اور نیکی میں اطاعت کی
        بیعت کی۔
  \item \textbf{عقبہ ثانیہ} (\textenglish{622}ء) --- تہتر مرد اور دو خواتین
        نے نبی کریم \saw کی اسی طرح حفاظت کی بیعت کی جیسے اپنے اہل و عیال
        کی کرتے ہیں --- اسے \textbf{بیعتِ حرب} (جنگ کی بیعت) کہا جاتا ہے،
        کیونکہ اس میں جان کی حفاظت کا وعدہ شامل تھا۔
\end{itemize}

\medskip
\textbf{ہجرت کا واقعہ۔} قریش نے دار الندوہ میں نبی کریم \saw کے قتل کا
منصوبہ بنایا --- ہر قبیلے سے ایک نوجوان مل کر یکجا وار کرتا، تاکہ خون کا
بدلہ کسی ایک قبیلے سے نہ لیا جا سکے۔ اسی رات حضرت علیؓ نبی \saw کے بستر پر
سوئے، اور آپ \saw حضرت ابوبکرؓ کے ساتھ روانہ ہوئے۔ تین دن \textbf{غارِ ثور}
میں چھپے رہے --- عنکبوت کا جالا اور کبوتری کا گھونسلا اس کے منہ پر، جس نے
تلاش کرنے والوں کو بھٹکا دیا۔ عامر بن فہیرہؓ بکریاں چرا کر نشانِ قدم مٹاتے
رہے، اور عبد اللہ بن ابوبکرؓ خبریں پہنچاتے رہے۔

\medskip
راستے میں سراقہ بن مالک انعام کے لالچ میں پیچھا کرتا رہا مگر ہر بار اس کے
گھوڑے کے پاؤں ریت میں دھنس جاتے --- بعد میں وہ خود مسلمان ہوا۔ آٹھ ربیع
الاول کو قافلہ \textbf{قبا} پہنچا، جہاں \textbf{مسجدِ قبا} (اسلام کی پہلی
مسجد) کی بنیاد رکھی گئی، پھر مدینہ میں داخل ہو کر \textbf{مسجدِ نبوی} کی
تعمیر ہوئی۔ یہی ہجرت اسلامی کیلنڈر (ہجری تقویم) کی بنیاد بنی۔
\end{hal}

%---------------------------------------------------------------------
\begin{sawal}{سوال \textenglish{7} \ \textnormal{\small(خزاں \textenglish{2012}، سوال \textenglish{7} اور خزاں \textenglish{2017}، سوال \textenglish{4})}}
عقدِ مواخات اور میثاقِ مدینہ پر الگ الگ جامع نوٹ تحریر کریں۔
\end{sawal}

\begin{hal}
\textbf{عقدِ مواخات (بھائی چارہ)۔} ہجرت کے فوراً بعد نبی کریم \saw نے ہر
مہاجر کو ایک انصاری کے ساتھ \textbf{بھائی} قرار دیا --- مثلاً حضرت ابوبکرؓ
اور خارجہ بن زیدؓ، حضرت عمرؓ اور عتبان بن مالکؓ۔ انصار نے اپنے مال، گھر
اور کاروبار میں مہاجرین کو شریک کیا۔ ابتدا میں یہ بھائی چارہ \textbf{وراثت
کی بنیاد} بھی بنا، بعد میں سورۃ الانفال (\textenglish{75}) کے نزول پر یہ
حکم منسوخ ہو کر وراثت خونی رشتے تک محدود ہوئی، مگر اخوت کا جذبہ قائم رہا۔

\smallskip
\textbf{مقصد:} مہاجرین کی معاشی و جذباتی سہارے کی ضرورت پوری کرنا، اور
اوس و خزرج کی پرانی دشمنی کے بعد ایک نئی، دین کی بنیاد پر متحد امت تشکیل
دینا --- اسی کا ذکر آل عمران \textenglish{103} میں ہے: \ar{فَأَلَّفَ بَيْنَ
قُلُوبِكُمْ فَأَصْبَحْتُم بِنِعْمَتِهِ إِخْوَانًا}۔

\medskip
\textbf{میثاقِ مدینہ۔} مدینہ میں مسلمانوں کے ساتھ یہود کے تین بڑے قبیلے
(بنو قینقاع، بنو نضیر، بنو قریظہ) اور دیگر عرب قبائل بھی آباد تھے۔ نبی کریم
\saw نے ایک تحریری معاہدہ مرتب کیا، جسے مؤرخین \textbf{دنیا کا پہلا تحریری
آئین} کہتے ہیں۔ اہم نکات:

\begin{itemize}\setlength{\itemsep}{1pt}
  \item مسلمان --- خواہ کسی قبیلے کے ہوں --- ایک ہی امت شمار ہوں گے۔
  \item یہود اپنے دین اور مال میں آزاد ہوں گے، بشرطیکہ وہ ظلم یا غداری نہ
        کریں۔
  \item مدینہ پر بیرونی حملے کی صورت میں سب مل کر دفاع کریں گے۔
  \item باہمی تنازعات کا فیصلہ نبی کریم \saw کے پاس لے جایا جائے گا۔
\end{itemize}

\smallskip
\textbf{اہمیت۔} اس معاہدے نے مدینہ کو ایک منظم ریاست میں بدل دیا، جس میں
نبی کریم \saw بحیثیت سربراہِ مملکت اور آخری منصف تسلیم کیے گئے --- اسلامی
سیاسی نظام کی پہلی باقاعدہ دستاویز۔
\end{hal}

%---------------------------------------------------------------------
\begin{sawal}{سوال \textenglish{8} \ \textnormal{\small(خزاں \textenglish{2017}، سوال \textenglish{5})}}
غزوہ بدر کے اسباب، واقعات، نتائج اور اہمیت پر مفصل نوٹ لکھیں۔
\end{sawal}

\begin{hal}
\textbf{بنیادی معلومات۔} \textenglish{17} رمضان \textenglish{2}ھ
(\textenglish{624}ء)، مقام: بدر (مدینہ سے جنوب مغرب)۔ مسلمان
\textbf{تین سو تیرہ}، قریش تقریباً \textbf{ایک ہزار}۔

\medskip
\textbf{اسباب۔} ابو سفیان کا شام سے آنے والا تجارتی قافلہ مسلمانوں کے
راستے سے گزرا۔ نبی کریم \saw نے صحابہؓ کو قافلے کی نگرانی کے لیے نکلنے کا
حکم دیا؛ ابو سفیان نے راستہ بدل کر قافلہ بچا لیا، مگر قریش نے پھر بھی
مکہ سے ایک بڑا لشکر نکالا تاکہ بدر میں طاقت کا مظاہرہ کریں۔

\medskip
\textbf{واقعات۔} نبی کریم \saw نے مشورے کے بعد بدر کے کنویں پر پڑاؤ کیا،
تاکہ پانی پر کنٹرول حاصل ہو۔ جنگ کا آغاز روایتی مبارزت (فردی مقابلوں) سے
ہوا، جس میں حضرت حمزہؓ، حضرت علیؓ اور حضرت عبیدہؓ نے مکی سرداروں کو شکست
دی۔ عام لڑائی میں اللہ نے فرشتوں کی مدد بھیجی: \ar{إِذْ تَسْتَغِيثُونَ
رَبَّكُمْ فَاسْتَجَابَ لَكُمْ أَنِّي مُمِدُّكُم بِأَلْفٍ مِّنَ
الْمَلَائِكَةِ} (الانفال، \textenglish{9})۔

\medskip
\textbf{نتائج۔} قریش کے \textbf{ستر} سردار مارے گئے --- بشمول ابو جہل ---
اور \textbf{ستر} قید ہوئے۔ قیدیوں سے سلوک نرم رکھا گیا: جو فدیہ ادا نہ کر
سکے، انہیں دس دس مسلمان بچوں کو لکھنا پڑھانا سکھا کر آزاد کیا گیا۔

\medskip
\textbf{اہمیت۔} قرآن نے اسے \textbf{یومِ فرقان} (حق و باطل میں فیصلہ کا
دن) کہا (الانفال، \textenglish{41})۔ یہ اسلام کی پہلی بڑی فوجی فتح تھی،
جس نے مسلمانوں کے حوصلے اور عرب میں ان کی سیاسی حیثیت دونوں کو مستحکم
کیا۔
\end{hal}

%---------------------------------------------------------------------
\begin{sawal}{سوال \textenglish{9} \ \textnormal{\small(بہار \textenglish{2022}، سوال \textenglish{4})}}
تیسری ہجری کے اہم واقعات قلمبند کریں۔
\end{sawal}

\begin{hal}
تیسرا سالِ ہجری \textbf{غزوہ اُحد} کے واقعے کی وجہ سے نمایاں ہے --- اس کی
مکمل تفصیل (اسباب، واقعات، قرآنی تجزیہ) پہلے ہی گہرائی سے پڑھی جا چکی ہونی
چاہیے (سورۃ آل عمران، آیات \textenglish{121--179})۔ اس سال کے دیگر اہم
نکات:

\begin{itemize}\setlength{\itemsep}{1pt}
  \item \textbf{غزوہ اُحد} --- شوال \textenglish{3}ھ، جبلِ احد۔ ابتدائی
        فتح، تیر اندازوں کی غلطی سے پلٹا، ستر صحابہ شہید (بشمول حضرت
        حمزہؓ)، نبی کریم \saw خود زخمی ہوئے۔ سبق: حکم کی خلاف ورزی فتح کو
        شکست میں بدل دیتی ہے۔
  \item \textbf{شراب کی تدریجی حرمت کا آغاز} --- سورۃ النساء
        (\textenglish{43}) میں نشے کی حالت میں نماز کے قریب جانے سے منع
        کیا گیا، جو مکمل حرمت کی طرف پہلا قدم تھا۔
  \item \textbf{پردے کے احکام کا نزول} --- ازواجِ مطہراتؓ اور مومن عورتوں
        کے لیے حجاب کی ابتدائی ہدایات اسی دور میں آئیں۔
  \item \textbf{حضرت زینبؓ کا نکاح} --- نبی کریم \saw کی صاحبزادی حضرت
        زینبؓ اسی سال (بعض روایات میں دوسری ہجری) مکہ سے ہجرت کر کے مدینہ
        پہنچیں۔
\end{itemize}

\medskip
\textbf{یاد رکھنے کا نکتہ۔} پرچے میں یہ سوال عموماً اس لیے پوچھا جاتا ہے
تاکہ طالب علم پورے سال کے واقعات کو \textbf{ترتیب وار} یاد رکھے، نہ کہ صرف
احد کو الگ تھلگ سمجھے --- غزوہ اُحد کو مرکز بنا کر اسی کے آس پاس کے احکامات
کا ذکر کر دینا کافی ہے۔
\end{hal}

%---------------------------------------------------------------------
\begin{sawal}{سوال \textenglish{10} \ \textnormal{\small(خزاں \textenglish{2017}، سوال \textenglish{6} اور بہار \textenglish{2013}، سوال \textenglish{1--6})}}
غزوہ خیبر کے اسباب، واقعات اور نتائج پر جامع نوٹ لکھیں، نیز صلح حدیبیہ کا
اہلِ خیبر پر کیا اثر پڑا؟
\end{sawal}

\begin{hal}
\textbf{صلح حدیبیہ کا پس منظر (مختصراً)۔} \textenglish{6}ھ میں نبی کریم
\saw عمرے کے ارادے سے روانہ ہوئے؛ قریش نے روکا، اور دس سالہ صلح طے پائی ---
شرائط بظاہر مسلمانوں کے خلاف نظر آئیں (اس سال عمرہ نہ ہو، اگلے سال ہو؛
قریش کا کوئی نوجوان مسلمانوں کے پاس آئے تو واپس، مگر مسلمان کا قریش کے
پاس جائے تو واپس نہیں)، مگر قرآن نے اسے \textbf{فتح مبین} قرار دیا (الفتح،
\textenglish{1}) --- کیونکہ امن کے ماحول میں دعوت آزادانہ پھیلی، اور دو
سال میں اتنے لوگ مسلمان ہوئے جتنے پہلے انیس سالوں میں نہیں ہوئے تھے۔

\medskip
\textbf{اثر بر اہلِ خیبر۔} خیبر کے یہود مدینہ سے شمال میں ایک مضبوط قلعہ
بند بستی میں رہتے تھے اور مسلمانوں کے خلاف مسلسل سازشیں کر رہے تھے۔ صلح
حدیبیہ کے بعد مسلمان قریش کی طرف سے بے فکر ہو گئے اور \textbf{پوری توجہ
خیبر کی طرف} مبذول کر سکے --- یعنی صلح نے بالواسطہ خیبر کی فتح کی راہ
ہموار کی۔

\medskip
\textbf{غزوہ خیبر۔} محرم \textenglish{7}ھ۔ خیبر کے کئی مضبوط قلعے تھے،
جن میں سے \textbf{قموص} سب سے مستحکم تھا۔ ابتدا میں دو جھنڈے دیے گئے مگر
فتح نہ ہوئی؛ نبی کریم \saw نے فرمایا: ''کل جھنڈا اسے دوں گا جو اللہ اور
اس کے رسول سے محبت رکھتا ہے اور اللہ و رسول اسے محبوب رکھتے ہیں''
(\ar{لَأُعْطِيَنَّ الرَّايَةَ}) --- اگلے دن حضرت علیؓ کو جھنڈا دیا گیا، اور
ان کی قیادت میں قموص فتح ہوا۔

\medskip
\textbf{نتائج۔} تمام قلعے یکے بعد دیگرے فتح ہوئے۔ یہود کی زمینیں مسلمانوں
کی ملکیت میں آئیں، مگر بٹائی (نصف پیداوار) پر انہی یہود کو کاشت کے لیے
واپس دے دی گئیں۔ عرب میں سب سے مضبوط یہودی طاقت کا خاتمہ ہوا، اور مدینہ کی
شمالی سرحد محفوظ ہو گئی۔
\end{hal}

%---------------------------------------------------------------------
\begin{sawal}{سوال \textenglish{11} \ \textnormal{\small(بہار \textenglish{2013}، سوال \textenglish{5} اور بہار \textenglish{2022}، سوال \textenglish{5})}}
فتح مکہ اور غزوہ حنین کا جائزہ پیش کریں۔
\end{sawal}

\begin{hal}
\textbf{اسباب۔} صلح حدیبیہ کے تحت بنو خزاعہ مسلمانوں کے حلیف اور بنو بکر
قریش کے حلیف تھے۔ قریش نے خفیہ طور پر بنو بکر کی مدد کر کے بنو خزاعہ پر
حملہ کرایا --- یہ صلح کی صریح خلاف ورزی تھی۔

\medskip
\textbf{واقعات۔} نبی کریم \saw \textbf{دس ہزار} صحابہؓ کے ساتھ \textenglish{20}
رمضان \textenglish{8}ھ (\textenglish{630}ء) مکہ میں داخل ہوئے --- عملاً
کوئی مزاحمت نہ ہوئی۔ ابو سفیان نے پہلے ہی اسلام قبول کر لیا تھا۔ نبی کریم
\saw نے خانہ کعبہ کے \textbf{تین سو ساٹھ بت} توڑے اور عام معافی کا اعلان
کیا: \ar{لَا تَثْرِيبَ عَلَيْكُمُ الْيَوْمَ} --- آج تم پر کوئی گرفت نہیں۔
حتیٰ کہ سب سے سخت دشمنوں کو بھی معاف کر دیا گیا۔

\medskip
\textbf{غزوہ حنین۔} فتح مکہ کے فوراً بعد، شوال \textenglish{8}ھ میں، قبیلہ
ہوازن اور ثقیف نے مسلمانوں پر حملے کی تیاری کی۔ مسلمان \textbf{بارہ ہزار}
کی کثیر تعداد پر بھروسا کر بیٹھے --- ابتدا میں دشمن کے اچانک حملے سے صفیں
بکھر گئیں، مگر نبی کریم \saw ثابت قدم رہے اور اللہ نے مدد بھیجی:
\ar{وَيَوْمَ حُنَيْنٍ إِذْ أَعْجَبَتْكُمْ كَثْرَتُكُمْ فَلَمْ تُغْنِ
عَنكُمْ شَيْئًا} (التوبہ، \textenglish{25})۔

\medskip
\textbf{نتائج۔} مسلمانوں کی فیصلہ کن فتح ہوئی، بڑی تعداد میں مالِ غنیمت
ملا۔ سبق: کثرتِ تعداد پر نہیں، اللہ کی مدد پر بھروسا کرنا چاہیے۔ فتح مکہ
اور حنین دونوں نے پورے عرب پر اسلام کی حتمی برتری ثابت کر دی۔
\end{hal}

%---------------------------------------------------------------------
\begin{sawal}{سوال \textenglish{12} \ \textnormal{\small(بہار \textenglish{2022}، سوال \textenglish{8})}}
اسلامی ریاست کے دفاع کے لیے نبی کریم \saw نے جو اقدامات کیے، ان کی تفصیل
تحریر کریں۔
\end{sawal}

\begin{hal}
\textbf{(الف) خطوطِ تبلیغی --- سفارتی حکمتِ عملی۔} \textenglish{6--7}ھ میں
نبی کریم \saw نے دنیا کے بڑے بادشاہوں کو دعوتِ اسلام کے خطوط بھیجے:
\textbf{ہرقل} (روم)، \textbf{کسریٰ} (ایران)، \textbf{مقوقس} (مصر)،
\textbf{نجاشی} (حبشہ)، اور دیگر علاقائی حکمران۔ ہر خط ایک معتبر صحابیؓ کے
ہاتھ بھیجا گیا۔ یہ اقدام محض تبلیغ نہیں، بلکہ اسلامی ریاست کی
\textbf{بین الاقوامی حیثیت} کا اعلان بھی تھا۔

\medskip
\textbf{(ب) غزوۂ تبوک۔} رجب \textenglish{9}ھ --- خبر ملی کہ روم مدینہ پر
حملے کی تیاری کر رہا ہے۔ نبی کریم \saw \textbf{تیس ہزار} کے لشکر کے ساتھ
سخت گرمی میں دور دراز تبوک تک پہنچے (اسی لیے اسے \textbf{غزوۂ عسرہ}،
مشکلات کا غزوہ، کہا جاتا ہے)۔ کوئی جنگ نہ ہوئی، مگر اس پیش قدمی نے روم پر
اسلامی ریاست کی طاقت واضح کر دی، اور منافقین کا اصل چہرہ بھی بے نقاب ہوا
(جو پیچھے رہ گئے)۔

\medskip
\textbf{(ج) دفاعی تیاری کا مستقل نظام۔} مدینہ کی حفاظت کے لیے نگرانی کے
دستے، قبائل سے معاہدے، اور \textbf{غزوۂ خندق} (\textenglish{5}ھ) میں
حضرت سلمان فارسیؓ کے مشورے پر خندق کھودنا --- عرب میں پہلی بار استعمال
ہونے والی دفاعی حکمتِ عملی، جس نے احزاب کے متحدہ لشکر کو ناکام بنا دیا۔

\medskip
\textbf{(د) عامُ الوفود۔} فتح مکہ اور تبوک کے بعد، \textenglish{9}ھ کو
\textbf{''عام الوفود''} (وفود کا سال) کہا جاتا ہے --- پورے عرب سے قبائلی
وفود مدینہ آ کر اسلام قبول کرنے لگے، جو اسلامی ریاست کے استحکام کا واضح
ثبوت تھا۔
\end{hal}
